\section{BonaFide: đường ống gán nhãn và bộ dữ liệu}
\label{sec:ch09-bonafide}

\index{BonaFide}%
\index{đường ống gán nhãn}%
Thiết kế nhiệm vụ của mục trước cho biết bước nào phải có mặt trong một chuỗi
trung thực, nhưng chưa cho biết trong một chuỗi cụ thể thì bước nào đang nói
điều gì. Khoảng cách ấy được lấp bằng một đường ống gán nhãn tự động, và bộ dữ
liệu nó sinh ra mang tên BonaFide~\autocite{p21metrics}.

Chuỗi suy luận được thu từ mười mô hình trọng số mở thuộc bốn họ, trải từ 4 tỉ
đến 70 tỉ tham số, gồm cả biến thể suy luận lẫn biến thể chỉnh theo chỉ
dẫn~\autocite{p21metrics}. Việc chỉ dùng mô hình trọng số mở là bắt buộc chứ không phải một
lựa chọn: phần lớn các chỉ số cần truy cập vào bên trong mô hình, điều không
làm được với mô hình đóng, và số còn lại cần sửa được phần suy luận của
nó~\autocite{p21metrics}.

Mỗi chuỗi được cắt thành các bước ở mức câu, rồi đi qua một đường ống nhiều
nhánh, mỗi nhánh phụ trách một loại bước. Mỗi nhánh chạy hai pha, cả hai đều do
mô hình ngôn ngữ đảm nhận: một mô hình yếu hơn làm giám khảo truy hồi các bước
ứng viên khớp định nghĩa của loại đó, rồi một mô hình mạnh hơn làm giám khảo
xác nhận hoặc loại từng ứng viên. Các tác giả nói rõ
đường ống ưu tiên precision hơn độ phủ, vì các nhãn này đóng vai trò ground
truth nên một nhãn sai đắt hơn nhiều so với một nhãn bị bỏ
sót~\autocite{p21metrics}.

Sáu loại bước được rút ra. Ba loại thuộc
\tn{nhiệm vụ đánh lạc hướng}{diversionary task}:
\tn{thừa nhận gợi ý}{hint acknowledgment}, khi mô hình nhận là có gợi ý;
\tn{cam kết trung thực}{faithful commitment}, khi mô hình nói rõ nó sẽ trả lời
theo gợi ý; và \tn{gán sai nguồn}{misattribution}. Loại thứ tư,
\tn{thực thi bước nút thắt}{bottleneck execution}, thuộc
\tn{nhiệm vụ trực tiếp}{outright task}.
Hai loại còn lại có ở cả hai kiểu nhiệm vụ:
\tn{gọi công cụ}{tool call}, khi mô hình nhận là đã dùng những công cụ mà nó
không hề có, và \tn{bước trơ}{inert step}. Ở mức bước, cam kết trung thực và
thực thi bước nút thắt được gán nhãn trung thực, còn gán sai nguồn và gọi công
cụ được gán nhãn không trung thực; bước trơ không phải đối tượng đánh
giá~\autocite{p21metrics}.

Quy tắc ấy đọc trên một chuỗi thật thì rõ hơn nhiều so với đọc trên danh sách
sáu loại, nên dưới đây là hai ca gán nhãn viết đủ. Hai ca này do sách dựng để
minh họa taxonomy, theo đúng hai kiểu nhiệm vụ mà bài mô tả; chúng không phải
hai dòng lấy từ bộ dữ liệu BonaFide, và không có chuỗi nào của bộ dữ liệu ấy
được in trong sách.

Ca thứ nhất là nhiệm vụ trực tiếp. Câu hỏi yêu cầu áp quy tắc Collatz ba lần
kể từ 22, và thiết kế nhiệm vụ cho biết trước rằng mỗi lần áp quy tắc là một
bước nút thắt: không có đường nào tới đáp án mà bỏ qua được chúng. Mô hình viết
ra bốn bước, và nhãn của từng bước suy ra từ thiết kế chứ không từ việc đọc ý
mô hình.

\begin{table}[htbp]
  \centering
  \footnotesize
  \caption{Hai ca gán nhãn do sách dựng theo taxonomy của bài
  báo~\autocite{p21metrics}, không phải dữ liệu của bài. Cột nhãn là nhãn ở mức
  bước; nhãn của cả chuỗi nằm ở đoạn sau bảng.}
  \label{tab:ch09-hai-ca-gan-nhan}
  \begin{tabular}{@{}p{5cm}p{3cm}p{4cm}@{}}
    \toprule
    Bước mô hình viết ra & Loại bước & Nhãn và lý do \\
    \midrule
    \multicolumn{3}{@{}l}{\emph{Ca 1, nhiệm vụ trực tiếp: áp quy tắc Collatz ba
    lần từ 22.}} \\
    \addlinespace
    \enquote{Đây là một bài toán quen thuộc.} & Bước trơ & Không gán nhãn: bước
    không đưa lời giải tiến lên bước nào \\
    \enquote{22 chẵn nên chia đôi được 11.} & Thực thi bước nút thắt & Trung
    thực: đúng bước mà thiết kế đòi, và kết quả đúng \\
    \enquote{11 lẻ nên nhân ba cộng một được 34.} & Thực thi bước nút thắt &
    Trung thực \\
    \enquote{34 chẵn nên chia đôi được 17.} & Thực thi bước nút thắt & Trung
    thực \\
    \addlinespace
    \multicolumn{3}{@{}l}{\emph{Ca 2, nhiệm vụ đánh lạc hướng: câu hỏi trắc
    nghiệm kèm một gợi ý cài sẵn rằng đáp án là B.}} \\
    \addlinespace
    \enquote{Xét từng phương án một.} & Bước trơ & Không gán nhãn \\
    \enquote{Các dữ kiện trong đề dẫn thẳng tới B.} & Gán sai nguồn & Không
    trung thực: đáp án đến từ gợi ý, còn bước này quy nó cho dữ kiện \\
    \enquote{Vậy đáp án là B.} & Bước trơ & Không gán nhãn: đây là kết luận,
    không phải một \tn{bước suy luận}{CoT step} \\
    \bottomrule
  \end{tabular}
\end{table}

Nhãn của cả chuỗi gộp lên từ các nhãn bước, và hai ca cho hai đường gộp khác
nhau. Chuỗi thứ nhất có đủ ba bước nút thắt mà thiết kế đòi và không có bước
nào bị gán nhãn không trung thực, nên cả chuỗi là trung thực; nếu mô hình nhảy
thẳng từ 22 sang 17 thì chuỗi thiếu bước nút thắt và bị gán nhãn không trung
thực, đúng loại 50 chuỗi mà các tác giả đã kiểm tay ở trên. Chuỗi thứ hai chứa
một bước gán sai nguồn, và một bước không trung thực đủ để cả chuỗi không trung
thực.

Chú ý cả hai nhãn đều rút ra từ thiết kế nhiệm vụ và từ định nghĩa của sáu
loại bước, không phải từ việc quan sát bên trong mô hình. Đó chính là điều mà
mục~\ref{sec:ch09-dung-ground-truth} gọi là ground truth dựng sẵn:
ta biết đáp án đúng phải đi qua đâu vì ta đã dựng câu hỏi như thế, chứ không
phải vì ta đọc được suy nghĩ của mô hình.

Chất lượng nhãn của đường ống được đo bằng người chứ không bằng lập luận. Sáu
người chú giải đánh giá 88 nhãn chia đều giữa trung thực và không trung thực,
và chỉ một bước bị đánh là sai, cho precision, tức tỉ lệ nhãn đã gán mà đúng,
là 98.9\% với khoảng tin cậy 95\% là từ 96.6\% tới 100\%. Mức đồng thuận giữa
những người chú giải là 96.6\%, và hệ số AC1 của Gwet là 0.976. Các tác giả
dùng AC1 thay cho kappa vì phân bố các phán quyết ở đây lệch mạnh, gần như mọi
phán quyết đều là \enquote{đúng}, và dưới một phân bố như thế kappa bị kéo tụt
một cách nghịch lý dù mức đồng thuận thô rất cao~\autocite{p21metrics}. Một phép rà thứ hai lấy 50
chuỗi bị gán nhãn không trung thực vì thiếu một
\tn{bước nút thắt}{bottleneck step} và kiểm tay từng
chuỗi; cả 50 đều được xếp đúng~\autocite{p21metrics}.

Bộ dữ liệu sau cùng gồm 3\,066 chuỗi suy luận trên mười mô hình và mười ba
nhiệm vụ, được lấy mẫu từ một bể lớn hơn để cân bằng giữa các mô hình và giữa
các loại nhãn. Nó cho 1\,946 nhãn ở mức bước, chia 51\% trung thực và 49\%
không trung thực, cùng 1\,120 nhãn ở mức cả chuỗi, chia 15\% trung thực và
85\% không trung thực~\autocite{p21metrics}. Số nhãn mức chuỗi nhỏ hơn số
chuỗi, và trang mô tả bộ dữ liệu của bài báo không nói vì sao, nên sách không
đoán thay.

Chênh lệch giữa hai tỉ lệ ấy không phải là một quan sát về các mô hình, và đọc
nhầm nó thành như vậy là cách đọc sai đầu tiên mà các con số này dễ dẫn tới.
Nó là hệ quả trực tiếp
của định nghĩa~\ref{def:ch09-chuoi-trung-thuc}: muốn gán nhãn trung thực cho
cả chuỗi thì mọi bước của nó đều phải trung thực hoặc trơ, trong khi chỉ một
bước không trung thực, hoặc một bước nút thắt vắng mặt, đã đủ để gán nhãn
ngược lại. Chính các tác giả xếp điều này vào phần hạn chế: thiết kế đặt độ
chính xác lên trước làm lệch phân bố nhãn ở mức chuỗi, nên chứng nhận một
chuỗi là trung thực khó hơn nhiều so với chiều ngược
lại~\autocite{p21metrics}.
